\chapter{PORTUGOL}
\section{Estrutura do Código}
\begin{lstlisting}[caption={Clasee programa em Portugol}, label={lst:programa_portugol}]

programa
{
    funcao inicio()
    {
        inteiro opcaoExec = 0
        inteiro opcaoMenu
        inteiro editarOp

        cadeia nomeCliente, sobrenomeCliente, cpf, celular
        cadeia logradouro, bairro, numCasa, complemento, cep

        cadeia tela, materiais, prazo, obs
        real valor

        faca
        {
            escreva("Escolha uma operação:\n")
            escreva("1- Cadastro Cliente\n")
            escreva("2- Cadastro Pedido\n")
            escreva("3- Consulta Cliente\n")
            escreva("4- Consulta Pedido\n")

            leia(opcaoMenu)

            escolha(opcaoMenu)
            {
                caso 1:
                    escreva("Nome:\n")
                    leia(nomeCliente)

                    escreva("Sobrenome:\n")
                    leia(sobrenomeCliente)

                    escreva("CPF:\n")
                    leia(cpf)

                    escreva("Celular:\n")
                    leia(celular)

                    escreva("Logradouro:\n")
                    leia(logradouro)

                    escreva("Bairro:\n")
                    leia(bairro)

                    escreva("Número:\n")
                    leia(numCasa)

                    escreva("Complemento:\n")
                    leia(complemento)

                    escreva("CEP:\n")
                    leia(cep)

                    escreva("Cliente cadastrado!\n")
                    pare

                caso 2:
                    escreva("Tamanho da tela:\n")
                    leia(tela)

                    escreva("Materiais:\n")
                    leia(materiais)

                    escreva("Prazo:\n")
                    leia(prazo)

                    escreva("Valor:\n")
                    leia(valor)
                    valor = validarNum(valor)

                    escreva("Observações:\n")
                    leia(obs)

                    escreva("Pedido cadastrado!\n")
                    pare

                caso 3:
                    escreva("CPF do cliente:\n")
                    leia(cpf)

                    escreva("Editar pedido?\n1-Sim\n2-Não\n")
                    leia(editarOp)

                    escolha(editarOp)
                    {
                        caso 1:
                            escreva("Tamanho da tela:\n")
                            leia(tela)

                            escreva("Materiais:\n")
                            leia(materiais)

                            escreva("Prazo:\n")
                            leia(prazo)

                            escreva("Valor:\n")
                            leia(valor)
                            valor = validarNum(valor)

                            escreva("Observações:\n")
                            leia(obs)

                            escreva("Pedido editado!\n")
                            pare

                        caso 2:
                            pare

                        caso contrario:
                            escreva("Opção inválida\n")
                    }
                    pare

                caso 4:
                    escreva("Editar cliente\n")

                    escreva("Nome:\n")
                    leia(nomeCliente)

                    escreva("Sobrenome:\n")
                    leia(sobrenomeCliente)

                    escreva("CPF:\n")
                    leia(cpf)

                    escreva("Celular:\n")
                    leia(celular)

                    escreva("Logradouro:\n")
                    leia(logradouro)

                    escreva("Bairro:\n")
                    leia(bairro)

                    escreva("Número:\n")
                    leia(numCasa)

                    escreva("Complemento:\n")
                    leia(complemento)

                    escreva("CEP:\n")
                    leia(cep)

                    escreva("Cliente editado!\n")
                    pare

                caso contrario:
                    escreva("Opção inválida\n")
            }

            escreva("Deseja continuar?\n1-Sim\n2-Não\n")
            leia(opcaoMenu)

            escolha(opcaoMenu)
            {
                caso 1:
                    pare

                caso 2:
                    escreva("Programa encerrado.\n")
                    opcaoExec = 1
                    pare

                caso contrario:
                    escreva("Opção inválida\n")
                    opcaoExec = 1
            }

        } enquanto (opcaoExec == 0)
    }

    // ===== FUNÇÃO DE VALIDAÇÃO =====
    funcao real validarNum(real numero)
    {
        enquanto (numero < 0)
        {
            escreva("Número inválido, informe novamente:\n")
            leia(numero)
        }
        retorne numero
    }
}


\end{lstlisting}

\chapter{JAVA}
Esse capítulo apresenta a implementação do sistema de gerenciamento de 
comissões artísticas utilizando a linguagem de programação Java. 

\section{Estrutura do Código}

\begin{lstlisting}[caption={Classe Menu em Java}, label={lst:menu_java}]

    /*
 * ALTINO AVILLA, GUILHERME DINIZ, TACIARA ROSENO DOS SANTOS, VIVIAN PASSOS
 * ATELIER PESSOAL
 * 15/04/2026
 */

//TEXTOS GRANDES TROCARPRA NEXT LINE
// garantir nextline funcione
import java.util.Scanner;

public class Menu {
	// estudar split
//METODOS
	public static String ValidarNome(String nome) {
		if (nome.contains("altino")) {
			System.out.println("Nome inválido.");
			System.out.println("Insira novamente:");
			Scanner sc = new Scanner(System.in);
			nome = sc.nextLine();
			return ValidarNome(nome);
		}
		if (nome.matches("^[A-Za-z]+\\s[A-Za-z]+$")) {
			System.out.println("Insira nome completo:");
			Scanner sc = new Scanner(System.in);
			nome = sc.nextLine();
			return ValidarNome(nome);
		}
		return nome;
	}



	//METODOS FIM
	public static void main(String[] args) {
		Scanner sc = new Scanner(System.in);
		// MENU
		int opcaoExec = 0;
		int opcaoMenu;
		// var criar cliente
		String nomeCliente, sobrenomeCliente, cpf, celular, logradouro, bairro, numCasa, complemento, cep;
		// var criar pedido
		String tela, materiais, prazo, obs;
		double valor;
		// var editar pedido, reusar criar pedido(?)
		int editarOp;
		
		do {
		System.out.println("Escolha uma operação:");
		System.out.println("1- Cadastro Cliente");
		System.out.println("2- Cadastro Pedido");
		System.out.println("3- Consulta Cliente");
		System.out.println("4- Consulta Pedido");
		opcaoMenu = sc.nextInt();

		switch (opcaoMenu) {

//CRIAR CLIENTE
		case 1:
			// talvez mudar pra caixa de aviso
			System.out.println("Insira o nome do cliente:");
			nomeCliente = sc.next();
			System.out.println("Insira o sobrenome do cliente:");
			sobrenomeCliente = sc.next();
			System.out.println("Insira o CPF do cliente:");
			cpf = sc.nextLine();
			System.out.println("Insira o celular do cliente:");
			celular = sc.next();
			System.out.println("Insira o logradouro:");
			logradouro = sc.next();
			System.out.println("Insira o bairro:");
			bairro = sc.next();
			System.out.println("Insira o número da casa:");
			numCasa = sc.next();
			System.out.println("Insira o complemento:");
			complemento = sc.next();
			System.out.println("Insira o CEP:");
			cep = sc.next();

			// validar cada um
			// jogar pro banco
			System.out.println("Cliente Cadastrado");
			break;
// CRIAR PEDIDO
		case 2:
			// talvez mudar pra caixa de aviso
			System.out.println("Selecione o cliente:");
			// informação banco de dados
			System.out.println("Informe o tamanho da tela:");
			tela = sc.next();
			System.out.println("Informe os materias utilizados:");
			materiais = sc.next();
			System.out.println("Informe o prazo:");
			// pegar info data (?)
			prazo = sc.next();
			System.out.println("Informe o valor:");
			valor = sc.nextFloat();
			System.out.println("Informe as observações:");
			obs = sc.next();

			// validar cada um
			// jogar pro banco
			System.out.println("Pedido Cadastrado");
			break;
// CONSULTAR E EDITAR PEDIDO
		case 3:
			// talvez mudar pra caixa de aviso
			System.out.println("Selecione o CPF do cliente:");
			// informação banco de dados
			System.out.println("Selecione o pedido:");
			// informação banco de dados
// *MOSTRAR DADOS PEDIDO banco
			System.out.println("Editar pedido?");
			System.out.println("1-Sim, 2-Não");
			editarOp = sc.nextInt();
			switch (editarOp) {
			case 1:
				System.out.println("Informe o tamanho da tela:");
				tela = sc.next();
				System.out.println("Infrome os materias utilizados:");
				materiais = sc.next();
				System.out.println("Informe o prazo:");
				// pegar info data (?)
				prazo = sc.next();
				System.out.println("Informe o valor:");
				valor = sc.nextFloat();
				System.out.println("Informe as observações:");
				obs = sc.next();

				System.out.println("Pedido Editado");
				break;
			case 2:
				break;
			default:
				System.out.println("Opção inválida");
				return;
			}
			break;
//EDITAR CLIENTE			
		case 4:
			// talvez mudar pra caixa de aviso
			System.out.println("Selecione o cliente:");
			// informação banco de dados
			System.out.println("Insira o nome do cliente:");
			nomeCliente = sc.next();
			System.out.println("Insira o sobrenome do cliente:");
			sobrenomeCliente = sc.next();
			System.out.println("Insira o CPF do cliente:");
			cpf = sc.next();
			System.out.println("Insira o celular do cliente:");
			celular = sc.next();
			System.out.println("Insira o logradouro:");
			logradouro = sc.next();
			System.out.println("Insira o bairro:");
			bairro = sc.next();
			System.out.println("Insira o número da casa:");
			numCasa = sc.next();
			System.out.println("Insira o complemento:");
			complemento = sc.next();
			System.out.println("Insira o CEP:");
			cep = sc.next();

			// validar cada um
			// jogar pro banco
			System.out.println("Clente Editado");
			break;

		default:
			System.out.println("Opção inválida");
			return;
		}// fim switch menu

		System.out.println("Realizar outra operação?");
		System.out.println("1-Sim, 2-Não");
		opcaoMenu = sc.nextInt();
		switch (opcaoMenu) {

		case 1:
			break;
		case 2:
			System.out.println("Operação finalizada.");
			opcaoExec =1;
			break;
		default:
			System.out.println("Opção inválida");
			return;
		}
		}//fim "do"
		while (opcaoExec == 0);
	}
}


\end{lstlisting}